\section{Phase 5: Manual Review}


Automated tools catch \emph{known} vulnerability patterns. Manual review catches:

\begin{itemize}[nosep]
  \item \textbf{Business logic errors}: The code does what it says, but what it says is
    wrong. Example: a fee calculation that charges 0.3\% on the output instead of the
    input, giving arbitrageurs a profitable loop.
  \item \textbf{Economic attacks}: Flash loan attacks, oracle manipulation, sandwich
    attacks. These exploit \emph{interactions between contracts} that automated tools
    analyze in isolation.
  \item \textbf{Governance attacks}: Vote buying, proposal frontrunning, timelock bypass.
  \item \textbf{Upgrade risks}: Proxy storage collisions, initialization order dependencies.
\end{itemize}

\subsection{Manual Review Checklist}

\begin{enumerate}[nosep]
  \item Read every line of code. Not skim---read.
  \item For each external call: can the callee re-enter? What state is inconsistent at the call point?
  \item For each \texttt{msg.sender} check: can it be bypassed via delegatecall or proxy?
  \item For each arithmetic operation: what happens at the extremes (0, 1, type(uint256).max)?
  \item For each oracle read: what if the oracle returns stale data? Zero? Negative?
  \item For each token interaction: does it handle fee-on-transfer tokens? Rebasing tokens?
  \item For each access control: is it consistent across all functions? Can admin keys be rotated?
  \item For each storage write: is the order correct? Can it be frontrun?
\end{enumerate}

%% ============================================================
